\begin{proof}

\begin{enumerate}[label={$\langle1\rangle$\arabic*.},leftmargin=*]
\item The divisor of any fibre is nonzero and effective. The index argument C.1 therefore gives $H^2(Y,T_Y\otimes M)\ne0$ for every line bundle $M$. Suppose $R^2f_*(T_Y\otimes L)=0$. Twisting by a sufficiently large positive pullback from the base kills $H^1(R^1f_*)$, so Leray would give a line bundle $M=L\otimes f^*\OO(n)$ with $H^2(Y,T_Y\otimes M)=0$, a contradiction. Thus $R^2f_*(T_Y\otimes L)$ is nonzero. Its stalk has nonzero residue at some point, and top-degree base change gives the first nonvanishing in \eqref{cdp:eq:C.16}. Cartier-fibre adjunction and Serre duality give the second, with the exact line bundle displayed there.

\item For torus general fibre \SeriesRef{cdp:st:C2a} gives $\sum e(F_{\rm red})=e(Y)=2$. In the class of \ref{cdp:st:C4a}, \ref{cdp:st:A.4.6} makes every singular reduced contribution either zero or two, with two exactly for the rational hexagonal normalization; smooth finite torus quotients have zero Euler number. There is consequently exactly one positive fibre. Its smooth rational normalization has $\Pic^0=0$ and torsion-free integral second cohomology. Every ambient line bundle has torsion first Chern class, whose pullback is therefore zero. In particular both $\eta^*K_Y$ and $\eta^*L$ are trivial. Serre duality on the connected rational normalization gives \eqref{cdp:eq:C.17}.

\end{enumerate}
\end{proof}
